\documentclass[12pt]{article}
\usepackage[T1]{fontenc}
\usepackage[utf8]{inputenc}
\usepackage{lmodern}
\usepackage{textcomp}
\usepackage{enumitem}
\usepackage{geometry}
\geometry{margin=1in}
\begin{document}
\begin{center}
Answer Key - Health Sciences - Graduate Level Assessment 7 \
\small (Answers are ordered exactly as in the question paper.)
\end{center}

\section*{Multiple Choice}
\begin{enumerate}[label=\arabic*.]
\item \textbf{Answer:} E
\\ \textit{Options:} A) cartilage and respiratory epithelium.; B) Only respiratory epithelium; C) Only cartilage; D) Respiratory epithelium and cartilage; E) smooth muscle and respiratory epithelium.
\\ \textit{Note:} The walls of bronchioles are composed of smooth muscle that allows for regulation of airway diameter and respiratory epithelium that facilitates gas exchange and protects the airways from pathogens and debris.
\item \textbf{Answer:} A
\\ \textit{Options:} A) constantly activated, regardless of ATP levels.; B) more efficient than glycolysis.; C) slow compared with glycolysis.; D) not involved in energy production.; E) irreversible.
\\ \textit{Note:} The creatine kinase reaction is constantly activated to maintain a steady supply of ATP for cellular energy demands, independent of ATP levels, as it acts as a buffer to quickly regenerate ATP from ADP and phosphocreatine during periods of high energy consumption.
\item \textbf{Answer:} A
\\ \textit{Options:} A) Haemophilia A; B) Fragile X syndrome; C) Becker muscular dystrophy; D) Color blindness; E) Lesch-Nyhan syndrome
\\ \textit{Note:} Haemophilia A is an X-linked recessive disorder, meaning it is associated with genes on the X chromosome, and thus does show X-linked inheritance, making this choice incorrect; therefore, the question likely contains a typographical error or misrepresentation regarding the intended correct answer.
\item \textbf{Answer:} A
\\ \textit{Options:} A) Shot gun sequencing; B) Cultivation in human neural cells; C) Deep pyro sequencing (NGS); D) Monoclonal antibody techniques
\\ \textit{Note:} Shotgun sequencing is the correct answer because it is a high-throughput method used to rapidly sequence and assemble the genomes of new polyomaviruses from fragmented DNA, allowing for comprehensive analysis of their genetic material.
\item \textbf{Answer:} A
\\ \textit{Options:} A) Ammonia, urea and creatine.; B) Ammonia, hypoxanthine and creatinine.; C) Ammonia, creatine and uric acid.; D) Urea, creatine and uric acid.; E) Ammonia, urea and uric acid.
\\ \textit{Note:} During multiple sprint sports, the rapid breakdown of ATP leads to increased production of ammonia, urea, and creatine as byproducts of energy metabolism and protein catabolism, resulting in their elevated concentrations in the blood.
\end{enumerate}

\section*{True / False}
\begin{enumerate}[label=\arabic*.]
\setcounter{enumi}{5}
\item \textbf{Answer:} True
\\ \textit{Options:} A) True; B) False
\\ \textit{Note:} Older adults typically have extensive life experiences and accumulated knowledge, enabling them to perform well on tests of semantic memory, which assesses general knowledge and facts.
\item \textbf{Answer:} False
\\ \textit{Options:} A) True; B) False
\\ \textit{Note:} Biotin is primarily involved in carboxylation reactions, particularly in the metabolism of fatty acids and amino acids, and does not directly facilitate the conversion of glucose to glycogen, which is primarily regulated by insulin and other factors.
\item \textbf{Answer:} False
\\ \textit{Options:} A) True; B) False
\\ \textit{Note:} A universal vaccine for all types of flu has not yet been developed due to the virus's high mutation rate and the diversity of its strains, preventing a single vaccine from providing comprehensive protection.
\item \textbf{Answer:} False
\\ \textit{Options:} A) True; B) False
\\ \textit{Note:} Measuring blood pressure in an elevated arm can lead to inaccurate readings due to the effects of gravity, which can alter the blood flow dynamics and result in artificially lower or higher measurements compared to the true blood pressure levels when measured at heart level.
\item \textbf{Answer:} True
\\ \textit{Options:} A) True; B) False
\\ \textit{Note:} First cousins share, on average, 12.5\% of their genes because they inherit approximately 1/8 of their genetic material from each of their common grandparents, resulting in a total of 1/4 when considering both sides of the family tree.
\end{enumerate}

\section*{Long Form Response}
\begin{enumerate}[label=\arabic*.]
\setcounter{enumi}{10}
\item \textbf{Answer:} The relative risk ratio of 1.2 (95\% CI: 1.1 to 1.8) indicates a statistically significant overall increase in the risk of hip fractures among elderly women who exercise regularly. This suggests that, contrary to the common assumption that exercise generally promotes bone health, there may be specific factors or types of exercise that contribute to an increased risk of fracture in this population. The confidence interval does not include 1, reinforcing the significance of the finding. Therefore, while exercise is typically associated with benefits for bone density and overall health, this data points to the necessity of further investigation into the types of physical activity undertaken by elderly women and their specific impacts on fracture risk. It highlights the complexity of the relationship between exercise and bone health, suggesting that not all exercise may be equally beneficial for fracture prevention in older adults.
\\ \textit{Note:} correct because a relative risk ratio of 1.2 with a confidence interval that does not include 1 indicates a statistically significant increase in hip fracture risk among elderly women who exercise regularly, suggesting the need to reevaluate the types of exercise they engage in and their potential impact on bone health.
\item \textbf{Answer:} Dietary factors, particularly saturated fat, have been implicated in influencing the risk of lung cancer through several mechanisms. High intake of saturated fats can lead to obesity and chronic inflammation, both of which are established risk factors for various cancers, including lung cancer. Additionally, saturated fats may alter lipid metabolism and promote the production of pro-inflammatory cytokines, which can contribute to tumorigenesis. Furthermore, diets rich in saturated fats may affect the immune response and hinder the body's ability to combat cancer cell proliferation. Understanding these mechanisms underscores the importance of dietary choices in cancer prevention strategies.
\\ \textit{Note:} correct because it identifies the link between high saturated fat intake and obesity and chronic inflammation, both of which increase lung cancer risk, while also addressing how saturated fats can influence lipid metabolism and promote pro-inflammatory cytokines that contribute to cancer development.
\end{enumerate}

\vfill
\noindent\textit{End of Answer Key}
\end{document}